\section{颜色模型的能力边界与适用范围}\label{sec:boundary}

前五章分别处理了数据、统计、判决、独立测试与补全尝试。本章不引入新的实验，而是回到已经固定的结果，回答一个更基本的问题：这套以 Cr/Cb 阈值为核心的方法，究竟在什么意义上有效，又在什么意义上失效。本章用到的两项统计——F1 分布与精确率上限——均由已保存的逐图像计数和原始类别计数直接算出，没有重新预测，也没有调整参数；它们属于对固定结果的事后解释，不构成新一轮独立测试。

\subsection{方法的实际构成}\label{subsec:composition}

把最终流程拆开，只包含三个判断。第一，像素的 Cr、Cb 是否同时落在以训练均值为中心、以 $k$ 倍标准差为半宽的矩形内；第二，画面中面积最大的 8 连通前景块是否就是脸部；第三，不与图像边界连通的前景空洞是否属于脸部。三者分别对应颜色先验、单主体先验和闭合轮廓先验。补全只是把闭合先验做得更强，没有增加新的判别信息。

这三个先验都不涉及“脸长什么样”。方法没有使用形状、纹理、边缘或五官位置，并在颜色转换后主动丢弃了 Y 通道。因此，我们更愿意把它描述为“肤色检测加两个空间约束”，而不是“人脸分割”。后文对优劣的判断都建立在这一认识上：误差主要不来自阈值是否精细，而来自三个先验在具体图像上是否成立。

这种简化并非疏忽。丢弃 Y 是为了让判决只依赖色度，减少曝光变化的影响；采用最大连通域，是为了在不引入人脸检测器的前提下排除零散干扰；填孔则是为了在不使用五官标签的前提下恢复眼睛、眉毛和嘴部。三项选择都以“不引入额外模型”为代价，换取实现的简单与结果的可解释。这些代价在第~\ref{sec:test}~章与第~\ref{sec:completion}~章的失败样例中均有体现。

\subsection{三个工作点}

表~\ref{tab:points} 把三次判断的结果放在一起。纯阈值是颜色单独作用时的位置；初版加入区域与填孔；补全再加膨胀与腐蚀。

\begin{table}[H]
\centering\small
\setlength{\tabcolsep}{4pt}
\caption{三种方案在测试集上的位置与证据状态}\label{tab:points}
\begin{tabular}{lrrrrl}
\toprule
方案 & Precision & Recall & F1 & IoU & 证据状态 \\
\midrule
纯阈值 & 52.91 & 76.82 & 62.66 & 45.63 & 独立测试 \\
初版 & 66.43 & 75.55 & 70.69 & 54.67 & 独立测试 \\
补全后 & 63.63 & 79.55 & 70.70 & 54.68 & 开发后对比 \\
\bottomrule
\end{tabular}
\end{table}

初版把精确率从 52.91 拉到 66.43，召回几乎不动；补全把召回推到 79.55，精确率退回 63.63。初版与补全后落在同一条 F1 等值线附近；基线不在同一条线上，F1 低约 8 个百分点。图~\ref{fig:pr} 把这层关系画了出来。

\begin{figure}[H]
\centering
\begin{tikzpicture}
\begin{axis}[
    width=0.82\textwidth,
    height=0.60\textwidth,
    xlabel={召回率 $R$（\%）},
    ylabel={精确率 $P$（\%）},
    xmin=50, xmax=85,
    ymin=45, ymax=75,
    xtick={50,55,60,65,70,75,80,85},
    ytick={45,50,55,60,65,70,75},
    grid=major,
    grid style={draw=gray!20},
    legend style={at={(0.02,0.98)},anchor=north west,font=\small,draw=none,fill=none},
    tick label style={font=\small},
    label style={font=\small},
]
\addplot[domain=64:85,samples=80,dashed,draw=gray!70,thick,forget plot] {70.7*x/(2*x-70.7)};
\node[gray!80,font=\scriptsize,anchor=north west] at (axis cs:82,61.5) {$F1=0.707$};
\addplot[domain=70:85,samples=60,dashed,draw=gray!45,forget plot] {65*x/(2*x-65)};
\node[gray!60,font=\scriptsize,anchor=north west] at (axis cs:80,54.0) {$F1=0.65$};
\addplot[only marks,mark=square*,mark size=2.6pt,draw=blue!60!black,fill=blue!60!black] coordinates {(76.82,52.91)};
\addlegendentry{纯阈值}
\addplot[only marks,mark=*,mark size=2.8pt,draw=myteal,fill=myteal] coordinates {(75.55,66.43)};
\addlegendentry{初版（区域+填孔）}
\addplot[only marks,mark=o,mark size=3.2pt,draw=myteal,fill=white,mark options={line width=0.8pt}] coordinates {(79.55,63.63)};
\addlegendentry{补全后（开发后）}
\node[font=\scriptsize,anchor=south west] at (axis cs:75.55,66.43) {70.69};
\node[font=\scriptsize,anchor=north west] at (axis cs:79.55,63.63) {70.70};
\node[font=\scriptsize,anchor=north] at (axis cs:76.82,52.91) {62.66};
\end{axis}
\end{tikzpicture}
\caption{三种方案在测试集上的精确率—召回率位置。实心点为独立测试，空心点为开发后对比；虚线为 F1 等值线。}
\label{fig:pr}
\end{figure}

这张图说明，补全没有让模型“更准”，只是把判决边界向外推了一圈：被包进来的像素中只有 35.45\% 是脸。初版与补全后之间的选择，不是“哪个方法更好”，而是“愿意用多少精确率换多少召回”。

\subsection{有效性的来源}

从结果分布看，32 张测试图的初版 F1 并不呈连续分布。不低于 0.85 的有 11 张，其中 4 张超过 0.92；低于 0.2 的有 5 张；0.2 至 0.5 之间没有样本。补全后这个形状基本不变：11 张不低于 0.85，5 张低于 0.2，只有个别图像跨过区间边界。方法在多数图像上表现较好，在少数图像上几乎完全失败，中间状态很少。

在我们看来，这种两段式分布与第~\ref{subsec:composition}~节的结构假设一致。当脸部构成画面中占主导的肤色连通区域时，三个先验同时成立：颜色候选接近完整，最大连通域不会选错，填孔也能补回被包围的五官。当这一条件不成立时，错误会沿同一条链条被放大，而不是表现为分数的平缓下降。当前样本只有 32 张，这一形态也可能带有偶然性；但它至少提示，方法的成败更接近“条件是否满足”，而非“参数是否精细”。

在满足条件的图像上，方法的表现相当稳定。样本 00610 的原始候选为 80,328 像素，最大连通域保留 79,899 像素，填孔补回 9,490 个正确像素，最终 F1 为 0.930，召回率为 97.7\%。00585、00629、00538 等图的情况类似。这些图像的一个共同特征是正面或近正面、脸部占画面比例较大、光照较均匀，头发与背景的颜色和肤色可分。对这些图像，颜色候选本身已接近完整，后处理只需少量修补。

方法的第二个优势体现在对五官的补偿上。颜色规则本来不覆盖眼睛、眉毛和嘴部，但填孔把这一组的召回率从 32.6\% 提高到 68.3\%，补回 69,662 个正确像素；皮肤与鼻子组也从 66.4\% 提高到 76.3\%。这说明“脸是一个被肤色包围的连通区域”这一先验在实际图像中相当强。方法并不区分这些像素的语义类别，只是利用了它们被脸部轮廓包围这一拓扑事实。就本实验的目标而言，这一补偿是有效的；但它也解释了补偿的上限：只有被完整包围的区域才能被补回。

第三个优势是可解释与可复现。方法只有四个统计参数和一个整数 $k$，没有随机初始化，也没有在测试阶段重新估计任何量。每一次失败都可以定位到颜色候选、区域选择或填孔中的某一步：00609 的错误始于颜色候选，00718 的错误始于区域选择，00676 的误差则在填孔阶段被放大。这种可归因性使讨论不必停留在“效果不好”，而可以具体到“哪一条先验没有成立”。此外，方法不需要训练，单张图像只经历一次颜色变换和两次连通域操作，结构上比需要训练的方法轻。

\subsection{失效的来源}

方法的失效可以按来源分为四类：颜色、空间、语义，以及评价口径。前三类属于方法本身，第四类属于我们如何定义对错。表~\ref{tab:failure} 给出摘要。

\begin{table}[H]
\centering\small
\setlength{\tabcolsep}{4pt}
\caption{失效的四类来源}\label{tab:failure}
\begin{tabular}{p{1.5cm}p{4.0cm}p{6.0cm}p{3.0cm}}
\toprule
来源 & 不成立的假设 & 主要证据 & 调参能否解决 \\
\midrule
颜色 & 肤色落在训练矩形内 & 00609：候选 5,685 像素、56 块、召回 1.6\%；低亮度组 F1 下降 3.8 个百分点 & 部分缓解，误检会回升 \\
空间 & 存在占主导的脸部连通块 & 00718：最大块 58,716 像素，TP 仅 2,705；脖子占误检 23.4\% & 需更换区域评分 \\
语义 & 肤色像素都属于脸 & 脖子与双耳 37.4 万像素；精确率上限约 85\% & 不能 \\
评价 & 整脸真值与候选边界一致 & 逐图 Mask IoU 0.572，Boundary IoU 0.233 & 不能，属定义问题 \\
\bottomrule
\end{tabular}
\end{table}

颜色层。当图像存在明显偏色或照度不足时，真实肤色的色度会偏离训练分布，$k=1.25$ 的矩形随之把大量正确像素排除在外。00609 是典型：$k=1.25$ 下只剩 5,685 个候选像素，分成 56 个连通块，最大块 1,063 像素，最终召回率 1.6\%。亮度分组的结果与此一致：较低亮度组的初版 F1 为 52.8\%，比基线低 3.8 个百分点，正确像素减少约 9.6 万，漏检增加约 9.6 万。补全对这类图几乎没有作用：00609 的召回只从 1.59\% 到 1.78\%。需要说明，亮度分组是事后按真值区域平均 $Y$ 划分的，只有 11 张图，且亮度与场景、肤色、妆容混杂，因此这里只能说明低亮度条件下风险更高，不能把下降单独归因于亮度。

空间层。最大连通域假设画面中存在一个占主导地位的脸部候选。当与肤色相近的头发形成更大连通块时，这一假设直接失效。00718 中，棕色头发构成 58,716 像素的最大块，其中正确像素只有 2,705 个，误检 56,011 个；填孔随后又增加 110 个正确像素和 10,871 个误检像素，最终 F1 为 0.045。补全把召回从 5.12\% 推到 9.78\%，仍然无济于事；00578 中，补全甚至沿着头发继续扩张，新增 22,692 个误检像素。这个例子说明，面积只能表示区域大小，不能表示区域属于什么。类似地，脖子与脸部通常连成同一块，无法通过连通性切开；在最终误检中，头发与脖子合计占 84.0\%，其中脖子单独占 23.4\%。

语义层。这是方法的真正上限。耳朵和脖子本身也是皮肤，颜色规则没有理由排除它们。测试集中，脖子与双耳合计约 37.4 万像素，占全部有效像素的 4.5\%，而脸部正类约 215 万像素。即使颜色判决完全正确、不再产生任何头发或背景误检，只要结构仍是“最大连通域加填孔”，这部分像素仍会被计入前景，精确率的上限约为 85\%。当前初版的精确率为 66.4\%，与这一上限之间还有约 19 个百分点，说明头发与背景误检仍是主要矛盾；但即使把这些误检全部消除，纯颜色模型也无法达到更高的精确率。这一计算建立在“所有肤色像素都被判为正类”的假设上，实际方法可能因连通性排除部分耳朵，因此 85\% 是量级上的上限，而非严格界。它的意义在于指出：剩余误差中有一部分不是参数问题，而是“颜色”这一信息本身不包含“脸”这一语义。

评价口径。还有一层不确定性来自评价本身。整脸真值与颜色候选之间的错位是任务定义的一部分：眼睛、眉毛和嘴部被计入正类，却不参与颜色建模。边界质量也受到这一错位的影响。初版方案的逐图 Mask IoU 平均为 0.572，而 Boundary IoU 平均仅为 0.233；前者说明区域大体重合，后者说明轮廓并不可信。填孔产生的是由外围轮廓决定的平滑边界，而不是真实的发际线或下颌线；标注本身在这些位置也存在歧义。因此，本方法适合给出粗略区域，不适合用于需要精确边缘的场合。

\subsection{适用条件}

把上述结果整理成条件，可以得到一个相对清晰的适用范围。这些条件来自 32 张测试图的观察，属于经验总结，不是经过独立验证的判据。表~\ref{tab:applicability} 给出摘要。

\begin{table}[H]
\centering\small
\setlength{\tabcolsep}{4pt}
\caption{适用条件与实验依据}\label{tab:applicability}
\begin{tabular}{p{4.2cm}p{7.2cm}p{3.4cm}}
\toprule
场景条件 & 实验观察 & 判断 \\
\midrule
单人、近正面、脸部占比较大、光照均匀 & 11 张 F1 不低于 0.85，4 张超过 0.92 & 适用 \\
偏色或照度不足 & 00609 F1 为 0.031；低亮度组下降 3.8 个百分点 & 风险高，需预处理 \\
头发或背景与肤色接近 & 头发占误检 60.6\%；00718 F1 为 0.045 & 易失效 \\
脸与脖子相连 & 脖子占误检 23.4\% & 精度受限 \\
多人、侧脸、遮挡 & 本数据未覆盖 & 未验证，不宜外推 \\
需要精确边缘 & Boundary IoU 0.233 & 不适用 \\
\bottomrule
\end{tabular}
\end{table}

据此，我们把适用场景限定为：单人、正面或近正面、脸部在画面中占较大比例、光照较均匀且与背景有较明显色差的图像；任务应是粗粒度区域提取或前置筛选，而不是精确边缘。多人、侧脸、遮挡等情形在本数据中没有覆盖，我们不把它们列入适用范围，也不据此判断方法在那些情形下的表现。概括而言，这套方法是一个在受控条件下可用的肤色先验，其适用范围主要由拍摄条件决定，参数调节只能在其内部微调。

如果必须在初版与补全后之间选一个，判断依据不是 F1，而是任务更在意什么。以 F1 和边界为准，选初版；以覆盖完整为准，选补全后。两者的 F1 相同，差别在精确率与召回的分配；这也是第~\ref{sec:completion}~章反复出现的“工作点”一词的实际含义。

\subsection{与其他颜色方法的关系}

最后，把本方法放回颜色分割的谱系中，可以更清楚地看出它的位置。

基于 RGB 的规则更简单，但三个通道都混有亮度信息，阈值同时受光照强度与颜色影响；Cr/Cb 把亮度分离出去，代价是丢弃了 Y 中的明暗线索。本实验没有直接比较两者，00609 的结果只能说明色度通道同样会受白平衡影响，不能说明 RGB 规则一定更差。HSV 的色调对光照变化相对稳健，但皮肤通常落在低饱和区域，色调在低饱和时并不稳定。

训练分布本身还提供了另一条线索。Cr 与 Cb 的相关系数约为 $-0.76$，联合分布呈倾斜的椭圆，而当前判决使用的是轴对齐矩形。矩形的四个角会包含肤色分布之外的取值组合，而倾斜方向上的肤色又可能被矩形的边切掉。若改用马氏距离或椭圆边界，不必更换颜色空间，只需在现有均值、标准差之外增加一个相关系数，把“各通道独立”的假设换成“两通道联合”的假设。就我们看到的分布形态而言，这是改动最小、针对性最直接的一处。本实验没有实现这一方案，因此不能给出它可能带来的提升；这里只把它作为由数据直接支持的一个判断记录下来。

逐图自适应阈值是另一种常见思路。它可以根据每张图的色度分布重新确定范围，对偏色更稳健；但在本任务中，背景中的木质墙面、暖色衣物同样会形成大面积色度聚集，自适应阈值有可能把这些区域纳入候选，与最大连通域的假设冲突。是否采用，取决于场景中是否存在稳定的肤色参照，而不能只按“能否适应光照”来判断。

至于基于学习的分割方法，它们能够利用形状、纹理和语义信息，理论上可以越过语义层的限制。但这类方法需要标注数据与训练过程，其结果也不再能简单归因到“颜色范围是否合适”。对于本实验提出的问题——颜色本身能走多远——引入学习模型后，这一问题就不再被回答。因此，我们选择在颜色框架内把边界刻画清楚，而不是用更强的模型换取更高的分数。

\subsection{小结}

本章的结论可以概括为一句话：这套方法的精度上限由它看不见的信息决定，而不是由 $k$ 或核大小决定。颜色先验、单主体先验和闭合轮廓先验在受控条件下同时成立，方法就能达到 0.9 左右；其中任何一条不成立，误差就会沿流程被放大，而不是被后处理吸收。参数调整只能改变三条先验的松紧，不能补上它们缺少的形状与语义。就本实验而言，它的价值主要在于实现代价低、结果可解释、失败可归因，而不在精度上限。
